\chapter{Getting started}
\setwlchapter{getting-started}
If two electrons of opposite spin occupy a localized orbital,
removing the up electron must leave a down electron. Changing the
order in which the two creators were applied changes the sign of the
resulting ket, but cannot change its occupation. This distinction
between an amplitude and a probability is a good first test of an
operator calculation.

\section{Installation and a first session}
Install the release paclet with \fn{PacletInstall} in Mathematica
12.1 or later. The English help pages require Mathematica 14.0 or later.
The installation command, which contacts the release server, is:
\begin{lstlisting}[style=wl]
PacletInstall[
 "https://github.com/rokzitko/sneg/releases/latest/download/SNEG.paclet"]
\end{lstlisting}
Load the package \emph{before} entering a definition using its symbols.
The lower-case loader also works on existing SNEG notebooks. In the
Documentation Center open \texttt{paclet:SNEG/guide/SNEG}.
\begin{wl}
Needs["sneg`"];
$SnegVersion
\end{wl}
Choose a new operator name and declare its algebra. In the following
input \fn{d[]} means one spin-degenerate orbital, not an annihilator.
The latter needs both a type and a spin index.
\begin{wl}
snegfermionoperators[d];
u = d[AN, UP]; v = d[AN, DO];
TableForm[{{antikomutator[u, conj[u]], antikomutator[u, conj[v]]},
           {antikomutator[v, conj[u]], antikomutator[v, conj[v]]}},
  TableHeadings -> {{"up", "down"}, {"up", "down"}}]
\end{wl}
The unit diagonal expresses Pauli exclusion for a single spin-orbital;
the off-diagonal zeros express orthogonality. The rule says nothing
about the occupation of the orbital. That information belongs to a
state or density matrix.

\section{From an operator to a spectrum}
Give the level an energy $\epsilon$ and assign an extra cost $U$ when
both spins are present. The algebra itself is independent of their
values. A separate line declares the parameters real so that
conjugation later has the intended meaning.
\begin{wl}
snegrealconstants[\[Epsilon], U];
h = \[Epsilon] number[d[]] + U hubbard[d[]];
blocks = makematricesbzvc[h, qsbasisvc[{d[]}]];
TableForm[{#[[1, 1]], #[[1, 2]], #[[2, 1, 1]]} & /@ blocks,
  TableHeadings -> {None, {"Q", "Spin", "Energy"}}]
\end{wl}
There are four states but only three rows. The middle row represents
both projections of a spin $1/2$. The label $Q=N-1$ counts electrons
relative to one electron per site. If $U>0$, changing $\epsilon$ from
positive to negative first makes the spin doublet the ground state.
Only below $\epsilon=-U$ is it favourable to add the second electron.
Chapter~10 revisits these two addition thresholds as poles of a
Green's function.

\section{Reading operator output}
A dagger distinguishes creation from annihilation in the formatted
output. To inspect the actual expression, request \fn{InputForm};
use \fn{FullForm} to see every head and argument. Both are views.
For instance, the occupation polynomial represented by
\fn{hubbard[d[]]} is unchanged if its displayed product acquires a
minus sign from putting all creators on the left. Use \fn{snegtex}
when the result is needed as a formula in a paper. A useful final
step before export is to compare a matrix element, rather than a
particular ordering of symbols.
\casestudy{moment}{Local moment as an occupation measurement}
\begin{exercise}[Parity spectroscopy]\label{ex:first}
At $\epsilon=-U/3$, $U>0$, determine the addition and removal costs
from the ground multiplet. Which excited level would be invisible in
a calculation restricted to the same charge sector?
\end{exercise}
